\IfFileExists{stacks-project.cls}{%
\documentclass{stacks-project}
}{%
\documentclass{amsart}
}

% For dealing with references we use the comment environment
\usepackage{verbatim}
\newenvironment{reference}{\comment}{\endcomment}
%\newenvironment{reference}{}{}
\newenvironment{slogan}{\comment}{\endcomment}
\newenvironment{history}{\comment}{\endcomment}

% For commutative diagrams we use Xy-pic
\usepackage[all]{xy}

% We use 2cell for 2-commutative diagrams.
\xyoption{2cell}
\UseAllTwocells

% We use multicol for the list of chapters between chapters
\usepackage{multicol}

% This is generally recommended for better output
\usepackage{lmodern}
\usepackage[T1]{fontenc}

% For cross-file-references
\usepackage{xr-hyper}

% Package for hypertext links:
\usepackage{hyperref}

% For any local file, say "hello.tex" you want to link to please
% use \externaldocument[hello-]{hello}
\externaldocument[introduction-]{introduction}
\externaldocument[conventions-]{conventions}
\externaldocument[sets-]{sets}
\externaldocument[categories-]{categories}
\externaldocument[topology-]{topology}
\externaldocument[sheaves-]{sheaves}
\externaldocument[sites-]{sites}
\externaldocument[stacks-]{stacks}
\externaldocument[fields-]{fields}
\externaldocument[algebra-]{algebra}
\externaldocument[brauer-]{brauer}
\externaldocument[homology-]{homology}
\externaldocument[derived-]{derived}
\externaldocument[simplicial-]{simplicial}
\externaldocument[more-algebra-]{more-algebra}
\externaldocument[smoothing-]{smoothing}
\externaldocument[modules-]{modules}
\externaldocument[sites-modules-]{sites-modules}
\externaldocument[injectives-]{injectives}
\externaldocument[cohomology-]{cohomology}
\externaldocument[sites-cohomology-]{sites-cohomology}
\externaldocument[dga-]{dga}
\externaldocument[dpa-]{dpa}
\externaldocument[sdga-]{sdga}
\externaldocument[hypercovering-]{hypercovering}
\externaldocument[schemes-]{schemes}
\externaldocument[constructions-]{constructions}
\externaldocument[properties-]{properties}
\externaldocument[morphisms-]{morphisms}
\externaldocument[coherent-]{coherent}
\externaldocument[divisors-]{divisors}
\externaldocument[limits-]{limits}
\externaldocument[varieties-]{varieties}
\externaldocument[topologies-]{topologies}
\externaldocument[descent-]{descent}
\externaldocument[perfect-]{perfect}
\externaldocument[more-morphisms-]{more-morphisms}
\externaldocument[flat-]{flat}
\externaldocument[groupoids-]{groupoids}
\externaldocument[more-groupoids-]{more-groupoids}
\externaldocument[etale-]{etale}
\externaldocument[chow-]{chow}
\externaldocument[intersection-]{intersection}
\externaldocument[pic-]{pic}
\externaldocument[weil-]{weil}
\externaldocument[adequate-]{adequate}
\externaldocument[dualizing-]{dualizing}
\externaldocument[duality-]{duality}
\externaldocument[discriminant-]{discriminant}
\externaldocument[derham-]{derham}
\externaldocument[local-cohomology-]{local-cohomology}
\externaldocument[algebraization-]{algebraization}
\externaldocument[curves-]{curves}
\externaldocument[resolve-]{resolve}
\externaldocument[models-]{models}
\externaldocument[functors-]{functors}
\externaldocument[equiv-]{equiv}
\externaldocument[pione-]{pione}
\externaldocument[etale-cohomology-]{etale-cohomology}
\externaldocument[proetale-]{proetale}
\externaldocument[relative-cycles-]{relative-cycles}
\externaldocument[more-etale-]{more-etale}
\externaldocument[trace-]{trace}
\externaldocument[crystalline-]{crystalline}
\externaldocument[spaces-]{spaces}
\externaldocument[spaces-properties-]{spaces-properties}
\externaldocument[spaces-morphisms-]{spaces-morphisms}
\externaldocument[decent-spaces-]{decent-spaces}
\externaldocument[spaces-cohomology-]{spaces-cohomology}
\externaldocument[spaces-limits-]{spaces-limits}
\externaldocument[spaces-divisors-]{spaces-divisors}
\externaldocument[spaces-over-fields-]{spaces-over-fields}
\externaldocument[spaces-topologies-]{spaces-topologies}
\externaldocument[spaces-descent-]{spaces-descent}
\externaldocument[spaces-perfect-]{spaces-perfect}
\externaldocument[spaces-more-morphisms-]{spaces-more-morphisms}
\externaldocument[spaces-flat-]{spaces-flat}
\externaldocument[spaces-groupoids-]{spaces-groupoids}
\externaldocument[spaces-more-groupoids-]{spaces-more-groupoids}
\externaldocument[bootstrap-]{bootstrap}
\externaldocument[spaces-pushouts-]{spaces-pushouts}
\externaldocument[spaces-chow-]{spaces-chow}
\externaldocument[groupoids-quotients-]{groupoids-quotients}
\externaldocument[spaces-more-cohomology-]{spaces-more-cohomology}
\externaldocument[spaces-simplicial-]{spaces-simplicial}
\externaldocument[spaces-duality-]{spaces-duality}
\externaldocument[formal-spaces-]{formal-spaces}
\externaldocument[restricted-]{restricted}
\externaldocument[spaces-resolve-]{spaces-resolve}
\externaldocument[formal-defos-]{formal-defos}
\externaldocument[defos-]{defos}
\externaldocument[cotangent-]{cotangent}
\externaldocument[examples-defos-]{examples-defos}
\externaldocument[algebraic-]{algebraic}
\externaldocument[examples-stacks-]{examples-stacks}
\externaldocument[stacks-sheaves-]{stacks-sheaves}
\externaldocument[criteria-]{criteria}
\externaldocument[artin-]{artin}
\externaldocument[quot-]{quot}
\externaldocument[stacks-properties-]{stacks-properties}
\externaldocument[stacks-morphisms-]{stacks-morphisms}
\externaldocument[stacks-limits-]{stacks-limits}
\externaldocument[stacks-cohomology-]{stacks-cohomology}
\externaldocument[stacks-perfect-]{stacks-perfect}
\externaldocument[stacks-introduction-]{stacks-introduction}
\externaldocument[stacks-more-morphisms-]{stacks-more-morphisms}
\externaldocument[stacks-geometry-]{stacks-geometry}
\externaldocument[moduli-]{moduli}
\externaldocument[moduli-curves-]{moduli-curves}
\externaldocument[examples-]{examples}
\externaldocument[exercises-]{exercises}
\externaldocument[guide-]{guide}
\externaldocument[desirables-]{desirables}
\externaldocument[coding-]{coding}
\externaldocument[obsolete-]{obsolete}
\externaldocument[fdl-]{fdl}
\externaldocument[index-]{index}

% Theorem environments.
%
\theoremstyle{plain}
\newtheorem{theorem}[subsection]{Theorem}
\newtheorem{proposition}[subsection]{Proposition}
\newtheorem{lemma}[subsection]{Lemma}

\theoremstyle{definition}
\newtheorem{definition}[subsection]{Definition}
\newtheorem{example}[subsection]{Example}
\newtheorem{exercise}[subsection]{Exercise}
\newtheorem{situation}[subsection]{Situation}

\theoremstyle{remark}
\newtheorem{remark}[subsection]{Remark}
\newtheorem{remarks}[subsection]{Remarks}

\numberwithin{equation}{subsection}

% Macros
%
\def\lim{\mathop{\mathrm{lim}}\nolimits}
\def\colim{\mathop{\mathrm{colim}}\nolimits}
\def\Spec{\mathop{\mathrm{Spec}}}
\def\Hom{\mathop{\mathrm{Hom}}\nolimits}
\def\Ext{\mathop{\mathrm{Ext}}\nolimits}
\def\SheafHom{\mathop{\mathcal{H}\!\mathit{om}}\nolimits}
\def\SheafExt{\mathop{\mathcal{E}\!\mathit{xt}}\nolimits}
\def\Sch{\mathit{Sch}}
\def\Mor{\mathop{\mathrm{Mor}}\nolimits}
\def\Ob{\mathop{\mathrm{Ob}}\nolimits}
\def\Sh{\mathop{\mathit{Sh}}\nolimits}
\def\NL{\mathop{N\!L}\nolimits}
\def\CH{\mathop{\mathrm{CH}}\nolimits}
\def\proetale{{pro\text{-}\acute{e}tale}}
\def\etale{{\acute{e}tale}}
\def\QCoh{\mathit{QCoh}}
\def\Ker{\mathop{\mathrm{Ker}}}
\def\Im{\mathop{\mathrm{Im}}}
\def\Coker{\mathop{\mathrm{Coker}}}
\def\Coim{\mathop{\mathrm{Coim}}}

% Boxtimes
%
\DeclareMathSymbol{\boxtimes}{\mathbin}{AMSa}{"02}

%
% Macros for moduli stacks/spaces
%
\def\QCohstack{\mathcal{QC}\!\mathit{oh}}
\def\Cohstack{\mathcal{C}\!\mathit{oh}}
\def\Spacesstack{\mathcal{S}\!\mathit{paces}}
\def\Quotfunctor{\mathrm{Quot}}
\def\Hilbfunctor{\mathrm{Hilb}}
\def\Curvesstack{\mathcal{C}\!\mathit{urves}}
\def\Polarizedstack{\mathcal{P}\!\mathit{olarized}}
\def\Complexesstack{\mathcal{C}\!\mathit{omplexes}}
% \Pic is the operator that assigns to X its picard group, usage \Pic(X)
% \Picardstack_{X/B} denotes the Picard stack of X over B
% \Picardfunctor_{X/B} denotes the Picard functor of X over B
\def\Pic{\mathop{\mathrm{Pic}}\nolimits}
\def\Picardstack{\mathcal{P}\!\mathit{ic}}
\def\Picardfunctor{\mathrm{Pic}}
\def\Deformationcategory{\mathcal{D}\!\mathit{ef}}
\usepackage{cleveref}
\crefname{originalreference}{original reference}{original references}
\Crefname{originalreference}{Original reference}{Original references}
\makeatletter
\newlabel{categories-lemma-2-fibre-product-erase-factor}{{02XD}{source}{}{sourcepointer-1}{}}
\newlabel{categories-lemma-2-fibre-product-erase-factor@cref}{{[originalreference][0][]02XD}{[1][0][]source}{}{}{}}
\makeatother
\crefname{originalreference}{original reference}{original references}
\Crefname{originalreference}{Original reference}{Original references}
\makeatletter
\newlabel{categories-lemma-2-product-categories-over-C}{{0040}{source}{}{sourcepointer-2}{}}
\newlabel{categories-lemma-2-product-categories-over-C@cref}{{[originalreference][0][]0040}{[1][0][]source}{}{}{}}
\makeatother
\crefname{originalreference}{original reference}{original references}
\Crefname{originalreference}{Original reference}{Original references}
\makeatletter
\newlabel{categories-lemma-associativity-2-fibre-product}{{02XC}{source}{}{sourcepointer-3}{}}
\newlabel{categories-lemma-associativity-2-fibre-product@cref}{{[originalreference][0][]02XC}{[1][0][]source}{}{}{}}
\makeatother
\crefname{originalreference}{original reference}{original references}
\Crefname{originalreference}{Original reference}{Original references}
\makeatletter
\newlabel{categories-lemma-category-fibred-setoids-presheaves-products}{{04SD}{source}{}{sourcepointer-4}{}}
\newlabel{categories-lemma-category-fibred-setoids-presheaves-products@cref}{{[originalreference][0][]04SD}{[1][0][]source}{}{}{}}
\makeatother
\crefname{originalreference}{original reference}{original references}
\Crefname{originalreference}{Original reference}{Original references}
\makeatletter
\newlabel{categories-lemma-triple-2-fibre-product-pr02}{{04S7}{source}{}{sourcepointer-5}{}}
\newlabel{categories-lemma-triple-2-fibre-product-pr02@cref}{{[originalreference][0][]04S7}{[1][0][]source}{}{}{}}
\makeatother
\crefname{originalequation}{original equation}{original equations}
\Crefname{originalequation}{Original equation}{Original equations}
\makeatletter
\newlabel{equation-morphisms-spaces}{{04SW}{source}{}{sourcepointer-6}{}}
\newlabel{equation-morphisms-spaces@cref}{{[originalequation][0][]04SW}{[1][0][]source}{}{}{}}
\makeatother
\crefname{originallemma}{original lemma}{original lemmas}
\Crefname{originallemma}{Original lemma}{Original lemmas}
\makeatletter
\newlabel{lemma-base-change-by-space-representable-by-space}{{0300}{source}{}{sourcepointer-7}{}}
\newlabel{lemma-base-change-by-space-representable-by-space@cref}{{[originallemma][0][]0300}{[1][0][]source}{}{}{}}
\makeatother
\crefname{originallemma}{original lemma}{original lemmas}
\Crefname{originallemma}{Original lemma}{Original lemmas}
\makeatletter
\newlabel{lemma-characterize-representable-by-space}{{02ZX}{source}{}{sourcepointer-8}{}}
\newlabel{lemma-characterize-representable-by-space@cref}{{[originallemma][0][]02ZX}{[1][0][]source}{}{}{}}
\makeatother
\crefname{originallemma}{original lemma}{original lemmas}
\Crefname{originallemma}{Original lemma}{Original lemmas}
\makeatletter
\newlabel{lemma-representable-diagonal}{{045G}{source}{}{sourcepointer-9}{}}
\newlabel{lemma-representable-diagonal@cref}{{[originallemma][0][]045G}{[1][0][]source}{}{}{}}
\makeatother
\crefname{originalreference}{original reference}{original references}
\Crefname{originalreference}{Original reference}{Original references}
\makeatletter
\newlabel{spaces-groupoids-lemma-quotient-stack-2-coequalizer}{{044U}{source}{}{sourcepointer-10}{}}
\newlabel{spaces-groupoids-lemma-quotient-stack-2-coequalizer@cref}{{[originalreference][0][]044U}{[1][0][]source}{}{}{}}
\makeatother
\crefname{originalreference}{original reference}{original references}
\Crefname{originalreference}{Original reference}{Original references}
\makeatletter
\newlabel{spaces-groupoids-lemma-quotient-stack-morphisms}{{044V}{source}{}{sourcepointer-11}{}}
\newlabel{spaces-groupoids-lemma-quotient-stack-morphisms@cref}{{[originalreference][0][]044V}{[1][0][]source}{}{}{}}
\makeatother
\crefname{originalreference}{original reference}{original references}
\Crefname{originalreference}{Original reference}{Original references}
\makeatletter
\newlabel{stacks-lemma-isom-as-2-fibre-product}{{04SI}{source}{}{sourcepointer-12}{}}
\newlabel{stacks-lemma-isom-as-2-fibre-product@cref}{{[originalreference][0][]04SI}{[1][0][]source}{}{}{}}
\makeatother
\title{A 1-morphism from a space-representable stack to an algebraic stack yields a groupoid in algebraic spaces and a fully faithful canonical morphism from the quotient stack}
\author{The Stacks Project Authors}
\date{stacks-project:algebraic.tex:local-snapshot-20261012; GNU FDL 1.2 or later}
\begin{document}
\maketitle
\noindent\textit{Editorial note.} The groupoid-axiom check in (4) is 'clear from our description of objects' with the description spelled out in-span; carried as an explicit obligation. The source explicitly rejects a long abstract-nonsense transfer argument in favor of the quasi-inverse trick, which is executed in full. The setoid characterization input is not carded ('Omitted, but see Categories') and is declared external.\par\medskip
\noindent\textit{Allowed standard prerequisites (editorial).}\begin{itemize}
\item Diagonal criterion and base change. lemma-representable-diagonal (carded L26680) and lemma-base-change-by-space-representable-by-space (carded L26677). Give part (1).
\item 2-fibre product algebra. Categories, lemmas associativity-2-fibre-product, 2-fibre-product-erase-factor, triple-2-fibre-product-pr02 (cross-chapter, named). Give parts (2) and (3) with explicit object-level descriptions.
\item Setoid presheaf correspondence. lemma-characterize-representable-by-space (same chapter, NOT carded, declared external) and Categories, lemma category-fibred-setoids-presheaves-products (cross-chapter, named). Identify R x\_\{pr1,U,pr0\} R with R x\_\{s,U,t\} R for part (4).
\item Quotient stack machinery. Groupoids in Spaces, lemmas quotient-stack-2-coequalizer and quotient-stack-morphisms (cross-chapter, named). Produce f\_can and the description of Isom sheaves for parts (5) and (6).
\item Isom sheaf as 2-fibre product. Stacks, lemma isom-as-2-fibre-product (cross-chapter, named). Used in part (6).
\end{itemize}\medskip
\ifdefined\hypertarget\hypertarget{sourcepointer-1}{}\fi
\noindent\textit{Source reference.} Original reference 02XD: focused-seeds/geometry/S9-alg-map-space-into-stack/ref-tags.txt, source lines 5--5. Navigation pointer only: the counted proof cites this statement by \textbackslash{}ref as a standing input; the deduction the counted proof has to make is carried out inside the counted span, and no part of that deduction is delegated to this pointer.\par\medskip
\ifdefined\hypertarget\hypertarget{sourcepointer-2}{}\fi
\noindent\textit{Source reference.} Original reference 0040: focused-seeds/geometry/S9-alg-map-space-into-stack/ref-tags.txt, source lines 6--6. Navigation pointer only: the counted proof cites this statement by \textbackslash{}ref as a standing input; the deduction the counted proof has to make is carried out inside the counted span, and no part of that deduction is delegated to this pointer.\par\medskip
\ifdefined\hypertarget\hypertarget{sourcepointer-3}{}\fi
\noindent\textit{Source reference.} Original reference 02XC: focused-seeds/geometry/S9-alg-map-space-into-stack/ref-tags.txt, source lines 7--7. Navigation pointer only: the counted proof cites this statement by \textbackslash{}ref as a standing input; the deduction the counted proof has to make is carried out inside the counted span, and no part of that deduction is delegated to this pointer.\par\medskip
\ifdefined\hypertarget\hypertarget{sourcepointer-4}{}\fi
\noindent\textit{Source reference.} Original reference 04SD: focused-seeds/geometry/S9-alg-map-space-into-stack/ref-tags.txt, source lines 8--8. Navigation pointer only: the counted proof cites this statement by \textbackslash{}ref as a standing input; the deduction the counted proof has to make is carried out inside the counted span, and no part of that deduction is delegated to this pointer.\par\medskip
\ifdefined\hypertarget\hypertarget{sourcepointer-5}{}\fi
\noindent\textit{Source reference.} Original reference 04S7: focused-seeds/geometry/S9-alg-map-space-into-stack/ref-tags.txt, source lines 9--9. Navigation pointer only: the counted proof cites this statement by \textbackslash{}ref as a standing input; the deduction the counted proof has to make is carried out inside the counted span, and no part of that deduction is delegated to this pointer.\par\medskip
\ifdefined\hypertarget\hypertarget{sourcepointer-6}{}\fi
\noindent\textit{Source reference.} Original equation 04SW: sources/stacks-algebraic/pinned-original-source.tex, source lines 696--696. Navigation pointer only: the counted proof cites this statement by \textbackslash{}ref as a standing input; the deduction the counted proof has to make is carried out inside the counted span, and no part of that deduction is delegated to this pointer.\par\medskip
\ifdefined\hypertarget\hypertarget{sourcepointer-7}{}\fi
\noindent\textit{Source reference.} Original lemma 0300: sources/stacks-algebraic/pinned-original-source.tex, source lines 900--900. Navigation pointer only: the counted proof cites this statement by \textbackslash{}ref as a standing input; the deduction the counted proof has to make is carried out inside the counted span, and no part of that deduction is delegated to this pointer.\par\medskip
\ifdefined\hypertarget\hypertarget{sourcepointer-8}{}\fi
\noindent\textit{Source reference.} Original lemma 02ZX: sources/stacks-algebraic/pinned-original-source.tex, source lines 671--671. Navigation pointer only: the counted proof cites this statement by \textbackslash{}ref as a standing input; the deduction the counted proof has to make is carried out inside the counted span, and no part of that deduction is delegated to this pointer.\par\medskip
\ifdefined\hypertarget\hypertarget{sourcepointer-9}{}\fi
\noindent\textit{Source reference.} Original lemma 045G: sources/stacks-algebraic/pinned-original-source.tex, source lines 1337--1337. Navigation pointer only: the counted proof cites this statement by \textbackslash{}ref as a standing input; the deduction the counted proof has to make is carried out inside the counted span, and no part of that deduction is delegated to this pointer.\par\medskip
\ifdefined\hypertarget\hypertarget{sourcepointer-10}{}\fi
\noindent\textit{Source reference.} Original reference 044U: focused-seeds/geometry/S9-alg-map-space-into-stack/ref-tags.txt, source lines 10--10. Navigation pointer only: the counted proof cites this statement by \textbackslash{}ref as a standing input; the deduction the counted proof has to make is carried out inside the counted span, and no part of that deduction is delegated to this pointer.\par\medskip
\ifdefined\hypertarget\hypertarget{sourcepointer-11}{}\fi
\noindent\textit{Source reference.} Original reference 044V: focused-seeds/geometry/S9-alg-map-space-into-stack/ref-tags.txt, source lines 11--11. Navigation pointer only: the counted proof cites this statement by \textbackslash{}ref as a standing input; the deduction the counted proof has to make is carried out inside the counted span, and no part of that deduction is delegated to this pointer.\par\medskip
\ifdefined\hypertarget\hypertarget{sourcepointer-12}{}\fi
\noindent\textit{Source reference.} Original reference 04SI: focused-seeds/geometry/S9-alg-map-space-into-stack/ref-tags.txt, source lines 12--12. Navigation pointer only: the counted proof cites this statement by \textbackslash{}ref as a standing input; the deduction the counted proof has to make is carried out inside the counted span, and no part of that deduction is delegated to this pointer.\par\medskip
\begin{lemma}
\label{lemma-map-space-into-stack}
Let $S$ be a scheme contained in $\Sch_{fppf}$.
Let $\mathcal{X}$ be an algebraic stack over $S$.
Let $\mathcal{U}$ be an algebraic stack over $S$ which
is representable by an algebraic space.
Let $f : \mathcal{U} \to \mathcal{X}$ be a 1-morphism. Then
\begin{enumerate}
\item the $2$-fibre product
$\mathcal{R} = \mathcal{U} \times_{f, \mathcal{X}, f} \mathcal{U}$
is representable by an algebraic space,
\item there is a canonical equivalence
$$
\mathcal{U} \times_{f, \mathcal{X}, f} \mathcal{U}
\times_{f, \mathcal{X}, f} \mathcal{U} =
\mathcal{R} \times_{\text{pr}_1, \mathcal{U}, \text{pr}_0} \mathcal{R},
$$
\item the projection $\text{pr}_{02}$ induces via (2) a $1$-morphism
$$
\text{pr}_{02} :
\mathcal{R} \times_{\text{pr}_1, \mathcal{U}, \text{pr}_0} \mathcal{R}
\longrightarrow
\mathcal{R}
$$
\item let $U$, $R$ be the algebraic spaces representing
$\mathcal{U}, \mathcal{R}$ and $t, s : R \to U$ and
$c : R \times_{s, U, t} R \to R$ are the morphisms corresponding
to the $1$-morphisms
$\text{pr}_0, \text{pr}_1 : \mathcal{R} \to \mathcal{U}$
and
$\text{pr}_{02} :
\mathcal{R} \times_{\text{pr}_1, \mathcal{U}, \text{pr}_0} \mathcal{R} \to
\mathcal{R}$ above, then the quintuple $(U, R, s, t, c)$ is a groupoid in
algebraic spaces over $S$,
\item the morphism $f$ induces a canonical $1$-morphism
$f_{can} : [U/R] \to \mathcal{X}$
of stacks in groupoids over $(\Sch/S)_{fppf}$, and
\item the $1$-morphism $f_{can} : [U/R] \to \mathcal{X}$ is fully faithful.
\end{enumerate}
\end{lemma}

\begin{proof}
Proof of (1). By definition $\Delta_\mathcal{X}$ is representable
by algebraic spaces so
Lemma \ref{lemma-representable-diagonal}
applies to show that $\mathcal{U} \to \mathcal{X}$ is representable
by algebraic spaces. Hence the result follows from
Lemma \ref{lemma-base-change-by-space-representable-by-space}.

\medskip\noindent
Let $T$ be a scheme over $S$. By construction of the $2$-fibre product (see
Categories, Lemma \ref{categories-lemma-2-product-categories-over-C})
we see that the objects of the fibre category $\mathcal{R}_T$
are triples $(a, b, \alpha)$ where $a, b \in \Ob(\mathcal{U}_T)$
and $\alpha : f(a) \to f(b)$
is a morphism in the fibre category $\mathcal{X}_T$.

\medskip\noindent
Proof of (2). The equivalence comes from repeatedly applying
Categories, Lemmas \ref{categories-lemma-associativity-2-fibre-product} and
\ref{categories-lemma-2-fibre-product-erase-factor}.
Let us identify
$\mathcal{U} \times_\mathcal{X} \mathcal{U} \times_\mathcal{X} \mathcal{U}$
with
$(\mathcal{U} \times_\mathcal{X} \mathcal{U})
\times_\mathcal{X} \mathcal{U}$.
If $T$ is a scheme over $S$, then on fibre categories over $T$
this equivalence maps the object
$((a, b, \alpha), c, \beta)$ on the left hand side
to the object $((a, b, \alpha), (b, c, \beta))$ of the right hand side.

\medskip\noindent
Proof of (3). The $1$-morphism $\text{pr}_{02}$ is constructed in the proof of
Categories, Lemma \ref{categories-lemma-triple-2-fibre-product-pr02}.
In terms of the description of objects of the fibre category
above we see that $((a, b, \alpha), (b, c, \beta))$
maps to $(a, c, \beta \circ \alpha)$.

\medskip\noindent
Unfortunately, this is {\it not compatible} with our conventions on
groupoids where we always have $j = (t, s) : R \to U$, and we ``think''
of a $T$-valued point $r$ of $R$ as a morphism $r : s(r) \to t(r)$.
However, this does not affect the proof of (4), since the opposite of
a groupoid is a groupoid. But in the proof of (5) it is responsible
for the inverses in the displayed formula below.

\medskip\noindent
Proof of (4). Recall that the sheaf $U$ is isomorphic to the sheaf
$T \mapsto \Ob(\mathcal{U}_T)/\!\cong$, and
similarly for $R$, see
Lemma \ref{lemma-characterize-representable-by-space}.
It follows from
Categories,
Lemma \ref{categories-lemma-category-fibred-setoids-presheaves-products}
that this description is compatible with $2$-fibre products
so we get a similar matching of
$\mathcal{R} \times_{\text{pr}_1, \mathcal{U}, \text{pr}_0} \mathcal{R}$
and $R \times_{s, U, t} R$.
The morphisms $t, s : R \to U$ and $c : R \times_{s, U, t} R \to R$
we get from the general equality (\ref{equation-morphisms-spaces}).
Explicitly these maps are the transformations of functors that come
from letting $\text{pr}_0$, $\text{pr}_0$, $\text{pr}_{02}$
act on isomorphism classes of objects of fibre categories.
Hence to show that we obtain a groupoid in algebraic
spaces it suffices to show that for every scheme $T$ over $S$
the structure
$$
(\Ob(\mathcal{U}_T)/\!\cong,
\Ob(\mathcal{R}_T)/\!\cong,
\text{pr}_1, \text{pr}_0, \text{pr}_{02})
$$
is a groupoid which is clear from our description of objects of
$\mathcal{R}_T$ above.

\medskip\noindent
Proof of (5). We will eventually apply
Groupoids in Spaces,
Lemma \ref{spaces-groupoids-lemma-quotient-stack-2-coequalizer}
to obtain the functor $[U/R] \to \mathcal{X}$.
Consider the $1$-morphism $f : \mathcal{U} \to \mathcal{X}$.
We have a $2$-arrow $\tau : f \circ \text{pr}_1 \to f \circ \text{pr}_0$
by definition of $\mathcal{R}$ as the $2$-fibre product.
Namely, on an object $(a, b, \alpha)$ of $\mathcal{R}$ over $T$ it is
the map $\alpha^{-1} : b \to a$. We claim that
$$
\tau \circ \text{id}_{\text{pr}_{02}} =
(\tau \star \text{id}_{\text{pr}_0})
\circ
(\tau \star \text{id}_{\text{pr}_1}).
$$
This identity says that given an object
$((a, b, \alpha), (b, c, \beta))$ of
$\mathcal{R} \times_{\text{pr}_1, \mathcal{U}, \text{pr}_0} \mathcal{R}$
over $T$, then the composition of
$$
\xymatrix{
c \ar[r]^{\beta^{-1}} & b \ar[r]^{\alpha^{-1}} & a
}
$$
is the same as the arrow $(\beta \circ \alpha)^{-1} : a \to c$. This is
clearly true, hence the claim holds. In this way we see that all the
assumption of
Groupoids in Spaces,
Lemma \ref{spaces-groupoids-lemma-quotient-stack-2-coequalizer}
are satisfied for the structure
$(\mathcal{U}, \mathcal{R}, \text{pr}_0, \text{pr}_1, \text{pr}_{02})$
and the $1$-morphism $f$ and the $2$-morphism $\tau$.
Except, to apply the lemma we need to prove this holds
for the structure $(\mathcal{S}_U, \mathcal{S}_R, s, t, c)$
with suitable morphisms.

\medskip\noindent
Now there should be some general abstract nonsense
argument which transfer these data between the two, but it seems to
be quite long. Instead, we use the following trick.
Pick a quasi-inverse $j^{-1} : \mathcal{S}_U \to \mathcal{U}$
of the canonical equivalence $j : \mathcal{U} \to \mathcal{S}_U$ which comes
from $U(T) = \Ob(\mathcal{U}_T)/\!\!\cong$.
This just means that for every scheme $T/S$ and every
object $a \in \mathcal{U}_T$ we have picked out a particular
element of its isomorphism class, namely $j^{-1}(j(a))$.
Using $j^{-1}$ we may therefore see $\mathcal{S}_U$
as a subcategory of $\mathcal{U}$. Having chosen this subcategory
we can consider those objects $(a, b, \alpha)$ of $\mathcal{R}_T$
such that $a, b$ are objects of $(\mathcal{S}_U)_T$, i.e., such
that $j^{-1}(j(a)) = a$ and $j^{-1}(j(b)) = b$. Then it is clear that
this forms a subcategory of $\mathcal{R}$ which maps isomorphically
to $\mathcal{S}_R$ via the canonical equivalence
$\mathcal{R} \to \mathcal{S}_R$. Moreover, this is clearly compatible
with forming the $2$-fibre product
$\mathcal{R} \times_{\text{pr}_1, \mathcal{U}, \text{pr}_0} \mathcal{R}$.
Hence we see that we may simply restrict
$f$ to $\mathcal{S}_U$ and restrict $\tau$ to a transformation
between functors $\mathcal{S}_R \to \mathcal{X}$. Hence it is clear that
the displayed equality of
Groupoids in Spaces,
Lemma \ref{spaces-groupoids-lemma-quotient-stack-2-coequalizer}
holds since it holds even as an equality of transformations of functors
$\mathcal{R} \times_{\text{pr}_1, \mathcal{U}, \text{pr}_0} \mathcal{R}
\to \mathcal{X}$ before restricting to the subcategory
$\mathcal{S}_{R \times_{s, U, t} R}$.

\medskip\noindent
This proves that
Groupoids in Spaces,
Lemma \ref{spaces-groupoids-lemma-quotient-stack-2-coequalizer}
applies and we get our desired morphism of stacks
$f_{can} : [U/R] \to \mathcal{X}$. We briefly spell out how
$f_{can}$ is defined in this special case.
On an object $a$ of $\mathcal{S}_U$ over $T$
we have $f_{can}(a) = f(a)$, where we think of
$\mathcal{S}_U \subset \mathcal{U}$ by the chosen embedding above.
If $a, b$ are objects of $\mathcal{S}_U$ over $T$, then a morphism
$\varphi : a \to b$ in $[U/R]$ is by definition an object of the
form $\varphi = (b, a, \alpha)$ of $\mathcal{R}$ over $T$. (Note
switch.) And the rule in the proof of
Groupoids in Spaces,
Lemma \ref{spaces-groupoids-lemma-quotient-stack-2-coequalizer}
is that
\begin{equation}
\label{equation-on-morphisms}
f_{can}(\varphi) = \Big(f(a) \xrightarrow{\alpha^{-1}} f(b)\Big).
\end{equation}
Proof of (6). Both $[U/R]$ and $\mathcal{X}$ are stacks.
Hence given a scheme $T/S$ and objects $a, b$ of $[U/R]$
over $T$ we obtain a transformation of fppf sheaves
$$
\mathit{Isom}(a, b) \longrightarrow \mathit{Isom}(f_{can}(a), f_{can}(b))
$$
on $(\Sch/T)_{fppf}$. We have to show that this is an
isomorphism. We may work fppf locally on $T$, hence we may assume that
$a, b$ come from morphisms $a, b : T \to U$. By the embedding
$\mathcal{S}_U \subset \mathcal{U}$ above we may also think of $a, b$ as
objects of $\mathcal{U}$ over $T$. In
Groupoids in Spaces,
Lemma \ref{spaces-groupoids-lemma-quotient-stack-morphisms}
we have seen that the left hand sheaf is represented by the algebraic space
$$
R \times_{(t, s), U \times_S U, (b, a)} T
$$
over $T$. On the other hand, the right hand side is by
Stacks, Lemma \ref{stacks-lemma-isom-as-2-fibre-product}
equal to the sheaf associated to the following stack in setoids:
$$
\mathcal{X}
\times_{\mathcal{X} \times \mathcal{X}, (f \circ b, f \circ a)} T =
\mathcal{X}
\times_{\mathcal{X} \times \mathcal{X}, (f, f)}
(\mathcal{U} \times \mathcal{U})
\times_{\mathcal{U} \times \mathcal{U}, (b, a)} T =
\mathcal{R}
\times_{(\text{pr}_0, \text{pr}_1), \mathcal{U} \times \mathcal{U}, (b, a)} T
$$
which is representable by the fibre product displayed above.
At this point we have shown that the two $\mathit{Isom}$-sheaves
are isomorphic. Our $1$-morphism $f_{can} : [U/R] \to \mathcal{X}$ induces
this isomorphism on $\mathit{Isom}$-sheaves by
Equation (\ref{equation-on-morphisms}).
\end{proof}

\end{document}
